That is, the formal series $S(x,t,\hbar)$ given by~\eqref{eq:Riccati-solution} satisfies
the following differential equations which are equivalent to~\eqref{eq:Lax-scalar}:
\begin{gather}  \label{eq:Riccati-2}
 \hbar^{2}\left( \left(\frac{\p S}{\p x}\right)^{2} +
\frac{\p^{2} S}{\p x^{2}} \right)   =
\frac{\hbar}{x-q} \left( \hbar \frac{\p S}{\p x} - p \right)
+ \big(4x^{3}+2t x + p^{2} - 4q^{3}-2tq\big),   \\
\label{eq:dt-Riccati-2}
 \hbar\frac{\p S}{\p t}   =   \frac{1}{2(x-q)} \left(
\hbar \frac{\p S}{\p x} - p \right).
\end{gather}
\end{thm}

Thus, the {\em principal specialization}
(i.e., setting $z_{i} = z$ for all $i = 1,\dots,n$)
of the open free energies gives an isomonodromic WKB solution.
Theorem~\ref{conj:2} implies
\begin{gather} \label{eq:dS-is-P}
\frac{\p S}{\p x}(x,t,\hbar) = P^{(+)}(x,t,\hbar)
\end{gather}
holds (under a suitable choice of the branch of $\sqrt{x+2q_{0}}$).
The computational results in Section~\ref{subsection:WKB-for-scalar-Lax}
show that~\eqref{eq:dS-is-P}
holds up to~$\hbar^{4}$. A~full proof of Theorem~\ref{conj:2}
will be given in Section~\ref{section:proof} together
with that of Theorem~\ref{conj:1} below.

\begin{rem}
In the topological recursion~\eqref{eq:top-rec},
we take residues only at the ramif\/ication point $z = 0$.
Thus $W_{g,n}$'s def\/ined here
are dif\/ferent from those in~\cite{DM14-2};
in particular, our quantum curve~\eqref{eq:Lax-scalar}
has {\em infinitely many $\hbar$-corrections} as in~\eqref{eq:hbar-correction-in-quantization} and~\eqref{eq:hbar-correction-in-quantization2}
(but recovers the same spectral curve in the semi-classical limit).
\end{rem}

\begin{rem}
In Theorem~\ref{conj:2}, the choice of the lower end points
of the integral in~\eqref{eq:open-Fgn} is important.
Dif\/ferent choice also give a WKB solution of the f\/irst equation
in~\eqref{eq:Lax-scalar}, but it may not satisfy
the second equation in general.
\end{rem}



\subsection[Closed free energies and the $\tau$-function]{Closed free energies and the $\boldsymbol{\tau}$-function}

The other main result of this paper is giving
another proof of the known fact
about the relationship between the closed free energies
and the $\tau$-function of~$\PI$
(cf.~\cite{DGZJ,EO07}).


\begin{Def}[\protect{\cite[Def\/inition~4.3]{EO07}}]
Def\/ine the {\em closed free energy}
$F_{g} = F_{g}(t)$ for $g \ge 2$ by
\begin{gather*}
F_{g}(t) = \frac{1}{2\pi i (2-2g)}
\oint_{\gamma_{0}} \Phi(z) W_{g,1}(z),
\end{gather*}
where
\begin{gather*}
\Phi(z) = \int^{z}_{z_{0}} y(z)dx(z)
  = \frac{4}{5}z^{5} - 4q_{0} z^{3} +
\text{const}
\end{gather*}
and $z_{0}$ is a generic point.
Free energies $F_{0}$ and $F_{1}$ for $g = 0, 1$
are also def\/ined but in a dif\/ferent manner
(see \cite[Sections~4.2.2 and~4.2.3]{EO07} for the def\/inition).
\end{Def}

Note that $F_g$ def\/ined here is dif\/ferent from $F_{g,n}$
def\/ined in the previous subsection.
$F_g$'s are also called {\it symplectic invariants}
since they are invariant under symplectic transformations
of the spectral curve (see~\cite{EO07}).
Explicit computation shows that
\begin{gather*}
F_{0}(t)   =   -\frac{48 q_{0}^{5}}{5},
\qquad
F_{1}(t)   =   -\frac{1}{24} \log(-3q_{0}), \\
F_{2}(t)   =   \frac{7}{207360 q_{0}^{5}},
\qquad
F_{3}(t)   =   \frac{245}{429981696 q_{0}^{10}}.
\end{gather*}

\begin{thm} [\protect{\cite{DGZJ} and \cite[Section~10.6]{EO07}}] \label{conj:1}
The generating function of the free energy $F_{g}(t)$
gives a $\tau$-function of~$\PI$:
\begin{gather*}
\log\tau(t,\hbar) = \sum_{g=0}^{\infty} \hbar^{2g-2} F_{g}(t).
\end{gather*}
Namely,
\begin{gather} \label{eq:tau-conjecture-g}
\frac{dF_{g}(t)}{dt} = \sigma_{2g}(t).
\end{gather}
\end{thm}

The proof will be given in Section~\ref{section:proof}.
It is worth mentioning that the closed free energies specify one
particular $\tau$-function although there is an ambiguity
in Def\/inition~\ref{def:tau-function}.

\begin{prop}
For $g \ge 2$, we have
\begin{gather} \label{eq:normalization-of-Fg}
F_{g}(t) = \int^{t}_{\infty} \sigma_{2g}(t')dt'.
\end{gather}
\end{prop}

\begin{proof}
Let us describe the behavior of the  $W_{g,n}$'s when
$q_{0} \rightarrow \infty$ (i.e., $t \rightarrow \infty$).
When $q_{0}$ tends to $\infty$,
no singular point of the integrand in the right hand-side of~\eqref{eq:top-rec}
on the $z$-plane hits the integration cycle $\gamma_{0}$.
Thus, we can show that
\begin{gather*}
W_{g,n}(z_{1},\dots,z_{n}) = O\big(q_{0}^{-(2g-2+n)}\big)
\end{gather*}
for $2g-2+n \ge 0$. This implies that
\begin{gather*}
F_{g}(t) = O\big(q_{0}^{-(2g-2)}\big)
\end{gather*}
holds since $\Phi(z) \sim q_{0}$ as $q_{0} \rightarrow \infty$
(but we can verify that~$F_{g}$ for $g \ge 2$
has a stronger decay in the above explicit computations).
This completes the proof of~\eqref{eq:normalization-of-Fg}.
\end{proof}


\subsection{Asymptotics of Eynard--Orantin dif\/fernetials}
\label{section:asymptotics-Wgn}

The rest of this section will be devoted to
show some important properties of $W_{g,n}$ and $F_{g,n}$.
Firstly, we will describe the asymptotic behavior
of them near $z_{i} = \infty$.

\begin{lem} \label{lem:decay-Wgn-Fgn}\quad
\begin{itemize}\itemsep=0pt
\item[$(i)$]
For $2g-2+n \ge 0$, we have
\begin{gather} \label{eq:asympt-Wgn}
W_{g,n}(z_{1},\dots,z_{n}) =
\left(\frac{c_{g,n}}{z_{1}^{2}\cdots z_{n}^{2}}
+ O\big(z_{1}^{-4}\cdots z_{n}^{-4}\big) \right)dz_{1}\cdots dz_{n}
\end{gather}
as $z_{i} \rightarrow \infty$ for all $i=1,\dots,n$.
Here $c_{g,n} \in \bbC$ is a constant.
\item[$(ii)$] %
For $2g-2+n \ge 0$, we have
\begin{gather} \label{eq:asympt-Fgn}
F_{g,n}(z_{1},\dots,z_{n}) =
\frac{c'_{g,n}}{z_{1}\cdots z_{n}} +
O\big(z_{1}^{-3}\cdots z_{n}^{-3}\big), \qquad
c'_{g,n} \in \bbC,
\end{gather}
as $z_{i} \rightarrow \infty$ for all $i=1,\dots,n$.
\end{itemize}
\end{lem}

\begin{proof}
The f\/irst property \eqref{eq:asympt-Wgn} follows from the
analyticity of $W_{g,n}$ at $z_i = \infty$.
The second property~\eqref{eq:asympt-Fgn} follows
from~\eqref{eq:asympt-Wgn} immediately
because $F_{g,n}(z_{1},\dots,z_{n})$ doesn't have
a constant term due to the def\/inition \eqref{eq:open-Fgn}.
\end{proof}

As a corollary, the principal specialization
of open free energies satisf\/ies
\begin{gather} \label{eq:decay-of-Fgn-prin}
F_{g,n}(z,\dots,z) = O\big(z^{-n}\big)
\end{gather}
when $z \rightarrow \infty$.

\subsection{Variation of spectral curve}
\label{section:variation-formula}

There is a formula (for ``variation of spectral curves'')
that allows us to compute derivatives of~$W_{g,n}$ etc.\
with respect to the parameter~$t$.

\begin{thm} [cf.~\protect{\cite[Theorem~5.1]{EO07}}]\label{thm:variation}\quad
\begin{itemize}\itemsep=0pt
\item[$(i)$]
For $2g-2+n \ge 0$, we have
\begin{gather}
\frac{\p}{\p t} W_{g,n}(z(x_{1}), \dots, z(x_{n}))\nonumber\\
\qquad{}
= -2\Res_{x_{n+1}=\infty} z(x_{n+1})
W_{g,n+1}\bigl( z(x_{1}),\dots,z(x_{n}),z(x_{n+1}) \bigr).\label{eq:variation-Wgn}
\end{gather}
\item[$(ii)$]
For $g \ge 1$, we have
\begin{gather} \label{eq:dt-closed-Fg}
\frac{dF_{g}}{dt}(t)
 = - 2\Res_{x=\infty} z(x) W_{g,1}(z(x))
= -\Res_{z=\infty} z W_{g,1}(z).
\end{gather}
\item[$(iii)$]
For $2g-2+ n \ge 1$, we have
\begin{gather*} %\label{eq:dt-openFgn}
\frac{\p}{\p t} F_{g,n}(z(x_{1}),\dots,z(x_{n})) \\
\qquad{} =
- 2\Res_{x_{n+1}=\infty} z(x_{n+1})
d_{x_{n+1}} F_{g,n+1}(z(x_{1}),\dots,z(x_{n}),z(x_{n+1})),
\end{gather*}
or equivalently,
\begin{gather}
\frac{\p}{\p t} F_{g,n}(z(x_{1}),\dots,z(x_{n})) \nonumber\\
\qquad{}= \lim_{z_{n+1} \rightarrow \infty} \left.
\left( z_{n+1}^{2}
\frac{\p}{\p z_{n+1}}F_{g,n+1}(z_{1},\dots,z_{n},z_{n+1})\right)
\right|_{(z_{1},\dots,z_{n})=(z(x_{1}), \dots, z(x_{n}))}. \label{eq:variation-openFgn-2}
\end{gather}
\end{itemize}
\end{thm}

\begin{proof}
Set $\Lambda(z) := z$.
Then, we can check $\Lambda(z)$ satisf\/ies
the required condition
\begin{gather*} % \begin{gather}
\Res_{z=\infty} \left( \Lambda(z) W_{0,2}(z,z_{1}) \right) = - dz_{1}
= - \left(\frac{\p y}{\p t}(z_{1}) dx(z_{1}) -
\frac{\p x}{\p t}(z_{1}) dy(z_{1}) \right)
\end{gather*} % \end{equation}
to apply \cite[Theorem~5.1]{EO07}. Thus the claim~(i) and~(ii)
are proved. Integrating both hand-sides of
\eqref{eq:variation-Wgn}, we have~(iii).
\end{proof}

\subsection{Dif\/ferential recursion for open free energies}

Here we give a key theorem in the proof of our main results.
We have the following {\em differential recursion} which is
a modif\/ication of the one obtained in \cite{DM14, DM14-2}.

\begin{thm} \label{thm:diff-rec}
The open free energies for $2g-2+n \ge 2$ satisfy
the following equations
\begin{gather}
\frac{\p F_{g,n}}{\p z_{1}}(z_{1},\dots,z_{n}) =
\sum_{j=2}^{n} \frac{-2z_{j}}{z_{1}^{2}-z_{j}^{2}}
\biggl( \frac{1}{2y(z_{1})\frac{dx}{dz}(z_{1})}
\frac{\p F_{g,n-1}}{\p z_{1}}(z_{[\hat{j}]})
 - \frac{1}{2y(z_{j})\frac{dx}{dz}(z_{j})}
\frac{\p F_{g,n-1}}{\p z_{j}}(z_{[\hat{1}]}) \biggr)
\nonumber\\
\hphantom{\frac{\p F_{g,n}}{\p z_{1}}(z_{1},\dots,z_{n}) =}{}
- \frac{1}{2y(z_{1})\frac{dx}{dz}(z_{1})} \frac{\p^{2}}{\p u_{1} \p u_{2}}
\biggl( F_{g-1,n+1}(u_{1},u_{2},z_{[\hat{1}]}) \nonumber\\
\hphantom{\frac{\p F_{g,n}}{\p z_{1}}(z_{1},\dots,z_{n}) =}{}
+
\sum_{\substack{g_{1}+g_{2}=g \\ I \sqcup J = [\hat{1}]}}^{\rm stable}
F_{g_{1}, |I|+1}(u_{1},z_{I}) F_{g_{2}, |J|+1}(u_{2}, z_{J})
\biggr)\Biggr|_{u_{1}=u_{2}=z_{1}}
\nonumber\\
\hphantom{\frac{\p F_{g,n}}{\p z_{1}}(z_{1},\dots,z_{n}) =}{}
+ \frac{s}{\frac{dy}{dz}(s) \frac{dx}{dz}(s)(z_{1}^{2}-s^{2})}
\Biggl[ \sum_{j=2}^{n} \frac{-2 z_{j}}{z_{j}^{2}-s^{2}}
\frac{\p F_{g,n-1}}{\p z_{1}}(s,z_{[\hat{1}, \hat{j}]})
\nonumber\\
\hphantom{\frac{\p F_{g,n}}{\p z_{1}}(z_{1},\dots,z_{n}) =}{}
+ \frac{\p^{2}}{\p u_{1} \p u_{2}}
\biggl( F_{g-1,n+1}(u_{1},u_{2},z_{[\hat{1}]})\nonumber\\
\hphantom{\frac{\p F_{g,n}}{\p z_{1}}(z_{1},\dots,z_{n}) =}{}
 +
\sum_{\substack{g_{1}+g_{2}=g \\
I \sqcup J = [\hat{1}]}}^{\rm stable}
F_{g_{1}, |I|+1}(u_{1},z_{I}) F_{g_{2}, |J|+1}(u_{2}, z_{J})
\biggr)\Biggr|_{u_{1}=u_{2}=s} \Biggr].\label{eq:diff-rec}
\end{gather}
Here
%\begin{gather*}% \label{eq:sing}
$s = (3q_{0})^{1/2}$
%\end{gather*}
is a zero of~$y(z)$.
\end{thm}

\begin{proof}
This can be proved by a similar technique used in
\cite[Theorem~4.7]{DM14}, as follows. Integrating
the topological recursion relation~\eqref{eq:top-rec}
with respect to $z_{2}, \dots, z_{n}$, we have
\begin{gather}
\frac{\p}{\p z_{1}} F_{g,n}(z_{1},\dots,z_{n})  =
\frac{1}{2^{n-1}}\int_{\bar{z}_{2}}^{z_{2}} \cdots
\int_{\bar{z}_{n}}^{z_{n}} W_{g,n}(z_{1},\dots,z_{n})\nonumber \\
\hphantom{\frac{\p}{\p z_{1}} F_{g,n}(z_{1},\dots,z_{n})}{}
=
\frac{1}{2\pi i} \frac{1}{2^{n-1}} \oint_{\gamma_{0}}
K(z,z_{1}) R_{g,n}(z,z_{2},\dots,z_{n}), \label{eq:pre-diff-recursion}
\end{gather}
where
\begin{gather*}
R_{g,n}(z,z_{2},\dots,z_{n}) =
\sum_{j=2}^{n}
\Biggl[ \biggl(\int^{z_{j}}_{\bar{z}_{j}}W_{0,2}(z,z_{j}) \biggr)
\biggl(\int^{z_{[\hat{1}, \hat{j}]}}_{\bar{z}_{[\hat{1}, \hat{j}]}}
W_{g,n-1}(\bar{z}, z_{[\hat{1},\hat{j}]}) \biggr)\\
\hphantom{R_{g,n}(z,z_{2},\dots,z_{n}) =}{}
- \biggl(\int^{z_{j}}_{\bar{z}_{j}}W_{0,2}(\bar{z},z_{j}) \biggr)
\biggl(\int^{z_{[\hat{1}, \hat{j}]}}_{\bar{z}_{[\hat{1}, \hat{j}]}}
W_{g,n-1}(z, z_{[\hat{1},\hat{j}]})
\biggr) \Biggr]
\\
\hphantom{R_{g,n}(z,z_{2},\dots,z_{n}) =}{}
+ \int^{z_{[\hat{1}]}}_{\bar{z}_{[\hat{1}]}}
W_{g-1,n+1}(z,\bar{z},z_{[\hat{1}]})\\
\hphantom{R_{g,n}(z,z_{2},\dots,z_{n}) =}{}
 + \sum_{\substack{g_{1}+g_{2}=g \\
 I \sqcup J = [\hat{1}]}}^{\text{stable}}
\biggl( \int^{z_{I}}_{\bar{z}_{I}}
W_{g_{1}, |I|+1}(z,z_{I}) \biggr)
\biggl( \int^{z_{J}}_{\bar{z}_{J}}
W_{g_{2}, |J|+1}(\bar{z}, z_{J})
\biggr).
\end{gather*}
Here, for a set $L =\{\ell_{1}, \dots, \ell_{k} \}
\subset \{1,\dots,n \}$ of indices, we have used the notation
\begin{gather*}
\int^{z_{L}}_{\bar{z}_{L}} W_{g,n}(z_{1},\dots,z_{n}) :=
\int^{z_{\ell_{1}}}_{\bar{z}_{\ell_{1}}} \cdots
\int^{z_{\ell_{k}}}_{\bar{z}_{\ell_{k}}} W_{g,n}(z_{1},\dots,z_{n}).
\end{gather*}
On the $z$-plane, the integrand
$K(z,z_{1})R_{g,n}(z,z_{1},\dots,z_{n})$
in the right hand-side of~\eqref{eq:pre-diff-recursion}
has poles at
\begin{itemize}\itemsep=0pt
\item
at $z = z_{1}, \bar{z}_{1}$ which are poles of $K(z,z_{1})$,
\item
at $z = z_{2}, \dots, z_{n}, \bar{z}_{2}, \dots, \bar{z}_{n}$
which are poles of $W_{0,2}(z,z_{j})$ and $W_{0,2}(\bar{z},z_{j})$,
\item
at $z = s, \bar{s}$ which are poles of $K(z,z_{1})$,
\end{itemize}
and all of them are simple poles.
Then, the equalities
\begin{gather*}
\int^{z_{j}}_{\bar{z}_{j}} W_{0,2}(z,z_{j})   =
\left(\frac{1}{z-z_{j}} - \frac{1}{z-\bar{z}_{j}} \right)dz, \\
\frac{1}{2^{n-2}}
\int^{z_{[\hat{1}, \hat{j}]}}_{\bar{z}_{[\hat{1}, \hat{j}]}}
W_{g,n}(z,z_{[\hat{1}, \hat{j}]})   =
\frac{\p F_{g,n-1}}{\p z_{1}}(z,z_{[\hat{1}, \hat{j}]})
\end{gather*}
and the residue theorem show~\eqref{eq:diff-rec}.
\end{proof}

\begin{rem} \label{rem:difference-from-DM}
Note that the f\/irst two blocks in the right hand-side
of~\eqref{eq:diff-rec} coincide with that obtained
in~\cite{DM14, DM14-2}.
Unlike the case of \cite{DM14, DM14-2},
we need more terms arising from $z=s$ corresponding
to the singular point $(x,y)=(q_{0},0)$
of the spectral curve~\eqref{eq:sp-curve}
since it becomes a~(simple) pole of
the recursion kernel $K(z,z_{1})$.
It also worth mentioning that the
right hand-side of~\eqref{eq:diff-rec} doesn't have
singularity at $z_{j} = s$ for $j=1,\dots,n$.
\end{rem}

Using this dif\/ferential recursion, we can give
an alternative expression of~\eqref{eq:variation-openFgn-2}
as follows.

\begin{thm} \label{thm:variation-of-open-free-energy}
For $2g-2+n \ge 1$, the following holds:
\begin{gather} \label{eq:t-derivative-recursion}
\frac{\p }{\p t} F_{g,n}(z(x_{1}),\dots,z(x_{n}))
= E_{g,n}(z(x_{1}),\dots,z(x_{n})),
\end{gather}
where
\begin{gather}
E_{g,n}(z_{1},\dots,z_{n}) \nonumber\\
\qquad{}
:=\sum_{j=1}^{n}
\frac{2z_{j}}{2y(z_{j})\frac{dx}{dz}(z_{j})}
\frac{\p F_{g,n}}{\p z_{j}}(z_{1},\dots,z_{n}) +
\frac{s}{\frac{dy}{dz}(s) \frac{dx}{dz}(s)}
\sum_{j=1}^{n} \frac{-2 z_{j}}{z_{j}^{2}-s^{2}}
\frac{\p F_{g,n}}{\p u_{1}}(u_{1},z_{[\hat{j}]})
\Biggr|_{u_{1}=s}
\nonumber\\
\qquad\quad{}+ \frac{s}{\frac{dy}{dz}(s) \frac{dx}{dz}(s)}
\frac{\p^{2}}{\p u_{1} \p u_{2}}
\biggl( F_{g-1,n+2}(u_{1},u_{2},z_{1},\dots,z_{n})
\nonumber\\
\qquad\quad{}+ \sum_{\substack{g_{1}+g_{2}=g \\
I \sqcup J = \{1,\dots,n\}}}^{\rm stable}
F_{g_{1}, |I|+1}(u_{1},z_{I}) F_{g_{2}, |J|+1}(u_{2}, z_{J})
\biggr)\Biggr|_{u_{1}=u_{2}=s}.\label{eq:Egn}
\end{gather}
\end{thm}

\begin{proof}
The equality \eqref{eq:variation-openFgn-2} shows that
the left hand-side of~\eqref{eq:t-derivative-recursion}
coincides with
\begin{gather*}
\lim_{z_{n+1} \rightarrow \infty} z_{n+1}^{2}
\frac{\p}{\p z_{n+1}} F_{g,n+1}(z_{1},\dots,z_{n},z_{n+1})
\end{gather*}
after the substitution $z_{i} \mapsto z(x_{i})$ for $i=1,\dots,n$.
Then, the equality follows from
the asymptotic behavior \eqref{eq:asympt-Fgn} of $F_{g,n}$'s
and the above dif\/ferential recursion \eqref{eq:diff-rec}
for $2g-2+ (n+1) \ge 2$.
\end{proof}


\section{Proof of main theorems}
\label{section:proof}

\subsection{Strategy for the proof}

What we will show here is that the formal series $S(x,t,\hbar)$
def\/ined in \eqref{eq:Riccati-solution} satisf\/ies
the system of equations \eqref{eq:Riccati-2} and
\eqref{eq:dt-Riccati-2}. In addition, we will also prove
the equality \eqref{eq:tau-conjecture-g}.
These equalities will be proved by an induction as follows.

\begin{thm} \label{thm:induction-scheme}
Let $[\bullet]_{\hbar^{m}}$ be the coefficient of
$\hbar^{m}$ in a formal series $\bullet$ of $\hbar$.
For an even integer $k \ge 2$, assume that
\begin{gather}
\frac{\p S_{m}}{\p x}(x,t) = P_{m}(x,t)\qquad
\text{for} \quad m=0,\dots,k-1, \nonumber\\
\frac{\p S_{m}}{\p t}(x,t) = \left[ \frac{1}{2(x-q)}\left(
\hbar \frac{\p S}{\p x} - p \right) \right]_{\hbar^{m}} \qquad
\text{for} \quad m=0,\dots,k-1, \nonumber\\
\frac{d F_{g}}{d t}(t) = \sigma_{2g}(t)  \qquad
\text{for} \quad g = k/2
\label{eq:assume-A-1}
\end{gather}
holds. Here $P_{m}(x,t) = P^{(+)}_{m}(x,t)$
is the coefficient of $\hbar^{m-1}$ in the formal solution
$P^{(+)}(x,t,\hbar)$ of the Riccati equation~\eqref{eq:Riccati}
constructed in Section~{\rm \ref{subsection:WKB-for-scalar-Lax}},
and~$\sigma_{2g}$ is given in~\eqref{eq:parity-sigma}.
Then, we have
\begin{itemize}\itemsep=0pt
\item[$(A)$] %
The following equality holds for $m = k$ and $k+1$:
\begin{gather}
\label{eq:claim-A-1}
\left[
\hbar^{2}\left( \left(\frac{\p S}{\p x}\right)^{2} +
\frac{\p^{2} S}{\p x^{2}} \right)
\right]_{\hbar^{m}}
 =  \left[ 2\hbar^{2} \frac{\p S}{\p t} +
\big(4x^{3}+2t x + p^{2} - 4q^{3}-2tq\big) \right]_{\hbar^{m}}.
\end{gather}

\item[$(B)$]
The following equalities hold:
\begin{gather}
\label{eq:claim-A-3}
\frac{\p S_{k}}{\p x}(x,t)   =   P_{k}(x,t), \qquad
\frac{\p S_{k}}{\p t}(x,t) = \left[ \frac{1}{2(x-q)}\left(
\hbar \frac{\p S}{\p x} - p \right) \right]_{\hbar^{k}}, \\
\frac{\p S_{k+1}}{\p x}(x,t)   =   P_{k+1}(x,t), \qquad
\frac{\p S_{k+1}}{\p t}(x,t) = \left[ \frac{1}{2(x-q)}\left(
\hbar \frac{\p S}{\p x} - p \right) \right]_{\hbar^{k+1}}, \label{eq:claim-B-3}
\\
\label{eq:claim-B-4}
\frac{dF_{g}}{dt}(t)   =    \sigma_{2g}(t)
\qquad \text{for} \quad g=(k+2)/2.
\end{gather}
\end{itemize}
\end{thm}

It is obvious that our main theorems (Theorems~\ref{conj:2} and~\ref{conj:1})
follow from the statements
in~$(A)$ and~$(B)$. The rest of this section is devoted to
give a proof of~$(A)$ and~$(B)$.

\subsection[Proof of $(A)$]{Proof of $\boldsymbol{(A)}$}
\label{section:proof-of-A}

We emphasize that the results shown in Section~\ref{subsec:proof-A-1} below
are proved without using the assumption~\eqref{eq:assume-A-1}.
We also note that we only use the second equality
in assumption~\eqref{eq:assume-A-1} in Section~\ref{subsec:proof-A-2}
to prove~$(A)$.

\subsubsection{Computation of principal specializations}
\label{subsec:proof-A-1}

Def\/ine
\begin{gather}
G_{g,n}(z_{1},\dots,z_{n}) :=
\frac{\p F_{g,n}}{\p z_{1}}(z_{1},\dots,z_{n})-\sum_{j=2}^{n} \frac{-2z_{j}}{z_{1}^{2}-z_{j}^{2}}
\biggl( \frac{1}{2y(z_{1})\frac{dx}{dz}(z_{1})}
\frac{\p F_{g,n-1}}{\p z_{1}}(z_{[\hat{j}]})
\nonumber\\
\hphantom{G_{g,n}(z_{1},\dots,z_{n}) :=}{}
 - \frac{1}{2y(z_{j})\frac{dx}{dz}(z_{j})}
\frac{\p F_{g,n-1}}{\p z_{j}}(z_{[\hat{1}]}) \biggr)
\nonumber\\
\hphantom{G_{g,n}(z_{1},\dots,z_{n}) :=}{}
+ \frac{1}{2y(z_{1})\frac{dx}{dz}(z_{1})} \frac{\p^{2}}{\p u_{1} \p u_{2}}
\biggl( F_{g-1,n+1}(u_{1},u_{2},z_{[\hat{1}]})\nonumber\\
\hphantom{G_{g,n}(z_{1},\dots,z_{n}) :=}{}
+ \sum_{\substack{g_{1}+g_{2}=g \\ I \sqcup J = [\hat{1}]}}^{\rm stable}
F_{g_{1}, |I|+1}(u_{1},z_{I}) F_{g_{2}, |J|+1}(u_{2}, z_{J})
\biggr)\Biggr|_{u_{1}=u_{2}=z_{1}}.
\label{eq:Ggn}
\end{gather}

The technique developed in
\cite{DM14, DM14-2} enables us to show the following.

\begin{lem} [cf.~\protect{\cite[Theorem 6.5]{DM14}}]
\label{prop:partial-Riccati}
For $m \ge 2$, we have
\begin{gather} \label{eq:partial-Riccati}
\Biggl(\frac{2y(z)}{\frac{dx}{dz}(z)}
\sum_{\substack{2g-2+n=m \\ g\ge0, n\ge 1}}
\frac{G_{g,n}(z,\dots,z)}{(n-1)!} \Biggr)\Biggl|_{z=z(x)} =
\sum_{\substack{a+b=m+1 \\ a,b \ge 0}}
\frac{\p S_{a}}{\p x} \frac{\p S_{b}}{\p x} +
\frac{\p^{2} S_{m}}{\p x^{2}} - \frac{1}{x-q_{0}} \frac{\p S_{m}}{\p x}.
\end{gather}
\end{lem}

\begin{proof}
As is shown in \cite[Theorem 6.5]{DM14},
applying $\sum\limits_{2g-2+n=m} \frac{1}{(n-1)!}$
and the principal specia\-li\-zation to~\eqref{eq:Ggn}, we have
\begin{gather*}
\sum_{\substack{2g-2+n=m \\ g\ge0, n\ge 1}}
\frac{G_{g,n}(z,\dots,z)}{(n-1)!} = \frac{1}{2y(z)\frac{dx}{dz}(z)}
\Biggl( \sum_{\substack{a+b=m+1 \\ a,b \ge 2}}
\frac{\p S_{a}(x(z))}{\p z} \frac{\p S_{b}(x(z))}{\p z} +
\frac{\p^{2}S_{m}(x(z))}{\p z^{2}} \Biggr)
\\
\hphantom{\sum_{\substack{2g-2+n=m \\ g\ge0, n\ge 1}}
\frac{G_{g,n}(z,\dots,z)}{(n-1)!} =}{}
+ \frac{\p S_{m+1}(x(z))}{\p z}
+ \left\{ \frac{\p }{\p z}\left(
\frac{1}{2 y(z) \frac{dx}{dz}(z)} \right) \right\}
\frac{\p S_{m}(x(z))}{\p z}.
\end{gather*}
After the coordinate change $z = z(x)$, the right hand-side becomes
\begin{gather*}
\frac{\frac{dx}{dz}(z(x))}{2y(z(x))}\Biggl(
\sum_{\substack{a+b=m+1 \\ a,b \ge 2}}
\frac{\p S_{a}}{\p x} \frac{\p S_{b}}{\p x} +
\frac{\p^{2}S_{m}}{\p x^{2}} + 2y(z(x)) \frac{\p S_{m+1}}{\p x}
- \frac{1}{y(z(x))}\frac{\p y(z(x))}{\p x} \frac{\p S_{m}}{\p x} \Biggr).
\end{gather*}
Then, the desired equality \eqref{eq:partial-Riccati} follows
from the above equality and
\begin{gather*}
\frac{\p S_{0}}{\p x} = y(z(x)), \qquad
\frac{\p S_{1}}{\p x} =
-\frac{1}{2 y(z(x))} \frac{\p y(z(x))}{\p x} + \frac{1}{2(x-q_{0})}.\tag*{\qed}
\end{gather*}
\renewcommand{\qed}{}
\end{proof}

Note that the right hand-side of~\eqref{eq:partial-Riccati} coincides with
\begin{gather*}
\biggl[ \hbar^{2}\biggl( \left(\frac{\p S}{\p x}\right)^{2} +
\frac{\p^{2} S}{\p x^{2}} \biggr)  \biggr]_{\hbar^{m+1}}
- \frac{1}{x-q_{0}} \frac{\p S_{m}}{\p x}.
\end{gather*}
Thus, Lemma \ref{prop:partial-Riccati} relates
the principal specialization of~$G_{g,n}$
to the left hand-side of~\eqref{eq:claim-A-1}.
Next we also relate them to the right hand-side of~\eqref{eq:claim-A-1}.

\begin{lem} \label{lemma:sum-prin-G-and-E}
Let $E_{g,n}(z_{1},\dots,z_{n})$
be the functions defined by~\eqref{eq:Egn}.
Then, the following equality holds for $m \ge 2$
\begin{gather} \label{eq:sum-prin-G-and-E}
\sum_{\substack{2g-2+n=m \\ g \ge 0, n\ge 2 }}
\left(
\frac{2y(z)}{\frac{dx}{dz}(z)}
\frac{G_{g,n}(z,\dots,z)}{(n-1)!} -
\frac{2 E_{g,n-1}(z,\dots,z)}{(n-1)!}\right)\Biggr|_{z=z(x)}
 = -\frac{1}{x-q_{0}} \frac{\p S_{m}}{\p x}.
\end{gather}
\end{lem}

\begin{proof}
Theorem \ref{thm:diff-rec} shows that
\eqref{eq:Ggn} can also be written as
\begin{gather}
G_{g,n}(z_{1},\dots,z_{n}) =
\frac{s}{\frac{dy}{dz}(s) \frac{dx}{dz}(s)(z_{1}^{2}-s^{2})}
\sum_{j=2}^{n} \frac{-2 z_{j}}{z_{j}^{2}-s^{2}}
\frac{\p F_{g,n-1}}{\p z_{1}}(s,z_{[\hat{1}, \hat{j}]})\nonumber
\\
\hphantom{G_{g,n}(z_{1},\dots,z_{n}) =}{}
+ \frac{s}{\frac{dy}{dz}(s) \frac{dx}{dz}(s)(z_{1}^{2}-s^{2})}
\frac{\p^{2}}{\p u_{1} \p u_{2}}
\biggl( F_{g-1,n+1}(u_{1},u_{2},z_{[\hat{1}]}) \nonumber\\
\hphantom{G_{g,n}(z_{1},\dots,z_{n}) =}{}+
\sum_{\substack{g_{1}+g_{2}=g \\ I \sqcup J = [\hat{1}]}}^{\rm stable}
F_{g_{1}, |I|+1}(u_{1},z_{I}) F_{g_{2}, |J|+1}(u_{2}, z_{J})
\biggr)\Biggr|_{u_{1}=u_{2}=s}.\label{eq:Ggn-2}
\end{gather}
Taking the principal specialization of \eqref{eq:Ggn-2} and~\eqref{eq:Egn} with $n \mapsto n-1$, we have
\begin{gather}
\frac{2y(z)}{\frac{dx}{dz}(z)}
\frac{G_{g,n}(z,\dots,z)}{(n-1)!} -
\frac{2 E_{g,n-1}(z,\dots,z)}{(n-1)!} \nonumber\\
\qquad{} =
- \frac{4z}{2y(z)\frac{dx}{dz}(z)}
\left(\frac{1}{(n-1)!} \frac{\p}{\p z} F_{g,n-1}(z,\dots,z)\right)\label{eq:prin-G-and-E}
\end{gather}
for any $g \ge 0$ and $n \ge 2$ satisfying
$2g-2+n \ge 2$.
Then, summing up~\eqref{eq:prin-G-and-E}
for $g \ge 0$, $n \ge 2$ satisfying $2g-2+n = m$,
we obtain~\eqref{eq:sum-prin-G-and-E} after
the coordinate change $z = z(x)$.
\end{proof}

On the other hand, Theorem~\ref{thm:variation-of-open-free-energy} implies that
\begin{gather*} %\label{eq:tderivative-of-open-free-energy}
\sum_{\substack{2g-2+n=m \\ g \ge 0,\,  n\ge 2 }}
\frac{2E_{g,n-1}(z,\dots,z)}{(n-1)!} \Biggr|_{z=z(x)}
= 2\frac{\p}{\p t} S_{m}
= \left[ 2\hbar^{2} \frac{\p S}{\p t} \right]_{\hbar^{m+1}}
\end{gather*}
holds for $m \ge 2$.
Therefore we have the following.

\begin{lem} \label{lem:tderivative-and-partial-Riccati}
The equality
\begin{gather*}%\label{eq:tderivative-and-partial-Riccati}
\left[ \hbar^{2}\left( \left(\frac{\p S}{\p x}\right)^{2} +
\frac{\p^{2} S}{\p x^{2}} \right)  \right]_{\hbar^{m+1}} \\
\qquad{} =
\left[ 2\hbar^{2} \frac{\p S}{\p t} \right]_{\hbar^{m+1}} + \sum_{\substack{2g-2+n=m \\ g \ge 0,\, n\ge {1} }}
\frac{2y(z)}{\frac{dx}{dz}(z)}
\frac{G_{g,n}(z,\dots,z)}{(n-1)!} -
\sum_{\substack{2g-2+n=m \\ g \ge 0, n\ge {2} }}
\frac{2y(z)}{\frac{dx}{dz}(z)}
\frac{G_{g,n}(z,\dots,z)}{(n-1)!}
\end{gather*}
holds for $m \ge 2$.
\end{lem}


\subsubsection[Completion of the proof of $(A)$]{Completion of the proof of $\boldsymbol{(A)}$}
\label{subsec:proof-A-2}

Lemma \ref{lem:tderivative-and-partial-Riccati} implies
\begin{gather}
\left[ \hbar^{2}\left( \left(\frac{\p S}{\p x}\right)^{2}
+ \frac{\p^{2} S}{\p x^{2}} \right)  \right]_{\hbar^{m+1}}
= \left[ 2\hbar^{2} \frac{\p S}{\p t} \right]_{\hbar^{m+1}}
 \qquad \text{if $m$ is even},\nonumber\\
\left[ \hbar^{2}\left( \left(\frac{\p S}{\p x}\right)^{2}
+ \frac{\p^{2} S}{\p x^{2}} \right)  \right]_{\hbar^{m+1}}
= \left[ 2\hbar^{2} \frac{\p S}{\p t} \right]_{\hbar^{m+1}}
+ \frac{2y}{\frac{dx}{dz}} G_{(m+1)/2,1}
\qquad \text{if $m$ is odd.}
\label{eq:compare-coefficients}
\end{gather}
On the other hand, it follows from \eqref{eq:parity-sigma} that
\begin{gather*}
%\label{eq:compare-coefficients-2}
\bigl[4x^{3}+2tx + p^{2} -4q^{3} -2 tq \bigr]_{\hbar^{m+1}}
=
\begin{cases}
0 & \text{if $m$ is even}, \\
 %%%
2\sigma_{m+1} & \text{if $m$ is odd}.
\end{cases}
\end{gather*}
Therefore, under the assumption \eqref{eq:assume-A-1},
the desired equality~\eqref{eq:claim-A-1}
follows from~\eqref{eq:compare-coefficients}
and Lem\-ma~\ref{lem:Gg1} below.

\begin{lem} \label{lem:Gg1}
For $g \ge 2$, we have
\begin{gather} \label{eq:n-is-1-part-in-Ggn}
\frac{2y(z)}{\frac{dx}{dz}(z)} G_{g,1}(z) =
 2 \frac{dF_{g}}{dt}(t).
\end{gather}
\end{lem}

\begin{proof}
Firstly, we note that
\begin{gather} \label{eq:explicit-Gg1}
\frac{2y(z)}{\frac{dx}{dz}(z)} G_{g,1}(z) =
\frac{1}{4s^{2}}
\biggl(
\frac{\p^{2} F_{g-1,2}}{\p z_{1} \p z_{2}}(s,s) +
\sum_{\substack{g_{1}+g_{2} = g \\ g_{1}, g_{2} \ge 1}}
\frac{\p F_{g_{1}, 1}}{\p z_{1}}(s)
\frac{\p F_{g_{2}, 1}}{\p z_{2}}(s)
\biggr)
\end{gather}
holds. Using the dif\/fernetial recursion \eqref{eq:diff-rec}
for $n=1$, we have
\begin{gather*}
\frac{\p F_{g,1}}{\p z_{1}} (z) =
- \frac{1}{2y(z)\frac{dx}{dz}(z)}
\biggl(
\frac{\p^{2} F_{g-1,2}}{\p z_{1} \p z_{2}}(z,z) +
\sum_{\substack{g_{1}+g_{2} = g \\ g_{1}, g_{2} \ge 1}}
\frac{\p F_{g_{1}, 1}}{\p z_{1}}(z)
\frac{\p F_{g_{2}, 1}}{\p z_{1}}(z)
\biggr)
\\
\hphantom{\frac{\p F_{g,1}}{\p z_{1}} (z) =}{}
+ \frac{s}{\frac{dy}{dz}(s)\frac{dx}{dz}(s)(z^{2}-s^{2})}
\biggl(
\frac{\p^{2} F_{g-1,2}}{\p z_{1} \p z_{2}}(s,s) +
\sum_{\substack{g_{1}+g_{2} = g \\ g_{1}, g_{2} \ge 1}}
\frac{\p F_{g_{1}, 1}}{\p z_{1}}(s)
\frac{\p F_{g_{2}, 1}}{\p z_{1}}(s)
\biggr).
\end{gather*}
Then, Lemma~\ref{lem:decay-Wgn-Fgn}
implies that
\begin{gather*}
z\frac{\p F_{g,1}}{\p z_{1}} (z)dz
= z W_{g,1}(z)  = \frac{1}{8s^{2}} \biggl(
\frac{\p^{2} F_{g-1,2}}{\p z_{1} \p z_{2}}(s,s) +
\sum_{\substack{g_{1}+g_{2} = g \\ g_{1}, g_{2} \ge 1}}
\frac{\p F_{g_{1}, 1}}{\p z_{1}}(s)
\frac{\p F_{g_{2}, 1}}{\p z_{2}}(s)
\biggr) \frac{dz}{z} + O(1)
\end{gather*}
holds when $z \rightarrow \infty$.
Then the equality~\eqref{eq:n-is-1-part-in-Ggn}
follows from~\eqref{eq:dt-closed-Fg} and~\eqref{eq:explicit-Gg1}.
\end{proof}

\subsection[Proof of $(B)$]{Proof of $\boldsymbol{(B)}$}

One of the desired equality \eqref{eq:claim-A-3} is proved as follows.

\begin{lem} \label{prop:integral-for-Sm}
Under the assumption \eqref{eq:assume-A-1}, we have
\begin{gather}
\label{eq:claim-A-5}
\frac{\p S_{k}}{\p x}(x,t)  =   P_{k}(x,t), \\
\label{eq:claim-A-4}
S_{k}(x,t)   =   \int^{x}_{\infty} P_{k}(x',t)dx', \\
\label{eq:claim-A-2}
\frac{\p S_{k}}{\p t}(x,t)   =   \left[\frac{1}{2(x-q)}
\left(\hbar \frac{\p S}{\p x} - p \right) \right]_{\hbar^{k}}.
\end{gather}
\end{lem}

\begin{proof}
The equality \eqref{eq:claim-A-1} for $m=k$ and
the second equality in the assumption~\eqref{eq:assume-A-1}
imply
\begin{gather*}
\left[
\hbar^{2}\left( \left(\frac{\p S}{\p x}\right)^{2} +
\frac{\p^{2} S}{\p x^{2}} \right)
\right]_{\hbar^{k}}
= \left[ \frac{\hbar}{x-q}\left( \hbar \frac{\p S}{\p x} - p \right)
+ \big(4x^{3}+2t x + p^{2} - 4q^{3}-2tq\big) \right]_{\hbar^{k}}.
\end{gather*}
Thus ${\p S_{k}}/{\p x}$ and $P_{k}$ satisfy the
same equation~\eqref{eq:recursion-for-Pm} under our induction hypothesis.
Then the uniqueness of the solution of~\eqref{eq:recursion-for-Pm}
implies~\eqref{eq:claim-A-5}.

Since $S_{m}(x)$ for $m \ge 2$ decay when
$x \rightarrow \infty$
(cf.~\eqref{eq:decay-of-Fgn-prin}),
the equality~\eqref{eq:claim-A-4}
immediately follows from~\eqref{eq:claim-A-5}.
Then, the equality~\eqref{eq:dt-P} shows
\begin{gather*}
\frac{\p}{\p t} S_{k}(x,t)   =
\int^{x}_{\infty} \left[ \hbar \frac{\p}{\p t} P
\right]_{\hbar^{k}}dx
  =
\int^{x}_{\infty} \frac{\p}{\p x} \left[
\frac{\hbar P - p}{2(x-q)} \right]_{\hbar^{k}} dx
  =
\left[
\frac{1}{2(x-q)} \left( \hbar \frac{\p S}{\p x} - p \right)
\right]_{\hbar^{k}}.
%\label{eq:t-derivation-of-S-and-P-in-Section-4}
\end{gather*}
The last equality follows from the assumption~\eqref{eq:assume-A-1} and the fact that
$P_{m}(x,t)$'s decay when $x \rightarrow \infty$
for $m \ge 1$ (see Remark~\ref{rem:asymptotics}),
and
\begin{gather*}
\lim_{x\rightarrow\infty} \frac{P_{0}(x,t)}{(x-q_{0})^{2}} = 0.
\end{gather*}
Thus we have proved \eqref{eq:claim-A-2}.
\end{proof}

Since we have also already proved \eqref{eq:claim-A-1}
for $m = k+1$, we can prove \eqref{eq:claim-B-3}
by the same discussion as the proof of
Lemma \ref{prop:integral-for-Sm} above.
Then, f\/inally we obtain

\begin{lem} \label{prop:partial-taufunction-conjecture}
The equality \eqref{eq:claim-B-4} is true; namely, we have
\begin{gather}\label{eq:claim-B-5}
\frac{dF_{(k+2)/2}}{dt} = \sigma_{2k+2}.
\end{gather}
\end{lem}

\begin{proof}
It follows from
the equality \eqref{eq:compare-coefficients}
(for the odd number $m = k+1$) and Lemma~\ref{lem:Gg1} that
\begin{gather} \label{eq:S-rec-section3}
2 \frac{\p S_{0}}{\p x}\frac{\p S_{k+2}}{\p x} +
\sum_{\substack{a+b=k+2 \\ a, b \ge 1}}
\frac{\p S_{a}}{\p x}\frac{\p S_{b}}{\p x} +
\frac{\p^{2} S_{k+1}}{\p x^{2}}
- 2\frac{\p S_{k+1}}{\p t}
 = 2 \frac{d F_{(k+2)/2}}{dt}
\end{gather}
holds.
On the other hand, we know that $P_{k+2}$ satisf\/ies
\begin{gather} \label{eq:P-rec-section3}
2 P_{0}P_{k+2} +
\sum_{\substack{a+b=k+2 \\ a, b \ge 1}}
P_{a}P_{b} +
\frac{\p P_{k+1}}{\p x} - \left[ \frac{\hbar}{x-q}
\left( \hbar P-p \right) \right]_{\hbar^{k+2}}
 = 2 \sigma_{k+2}
\end{gather}
(cf.~\eqref{eq:recursion-for-Pm}).
Under our assumption, comparing \eqref{eq:S-rec-section3}
and \eqref{eq:P-rec-section3}, we have
\begin{gather} \label{eq:final-remark}
\frac{\p S_{0}}{\p x} \left( \frac{\p S_{k+2}}{\p x} - P_{k+2} \right)
 = \frac{d F_{(k+2)/2}}{dt} - \sigma_{k+2}.
\end{gather}
Note that the right hand-side doesn't depend on~$x$.
Then, thanks to the fact
\begin{gather*}
\frac{\p S_{0}}{\p x} \biggr|_{x=q_{0}}
= 0
\end{gather*}
and the holomorphicity of $S_{m}(x)$ and $P_{m}(x)$
at the double turning point $x = q_0$
(see Theorem~\ref{thm:Podd-and-t-derivative}),
we have the desired equality~\eqref{eq:claim-B-5}
by substituting $x = q_{0}$ into~\eqref{eq:final-remark}.
\end{proof}

This completes the proof of $(B)$ and Theorem~\ref{thm:induction-scheme}.
Thus we have proved Theorems~\ref{conj:2} and~\ref{conj:1}.
